\section{Anexo 1: Declaración de uso de IAG}
\label{sec:anexo_iag}

En cumplimiento de la normativa de la Universidad Intercontinental de la Empresa, se declara el uso de las siguientes herramientas de Inteligencia Artificial Generativa (IAG) durante la elaboración de este Trabajo Fin de Grado:

\subsection*{Herramientas utilizadas}

\begin{enumerate}
  \item \textbf{Claude (Anthropic) --- familia Claude 4.x (\texttt{claude-opus-4-7}, \texttt{claude-sonnet-4-6}) vía Claude Code en VSCode.}
  \begin{itemize}
    \item \textbf{Uso}: Asistencia en la generación y refactorización de código Rust (tipos de transacción, consenso PoA, EVM, pool, red P2P), redacción y revisión técnica de la propuesta de EIP, estructuración y mejora de redacción de la memoria, depuración de errores de compilación, generación de diagramas TikZ y verificación matemática de los ejemplos trabajados.
    \item \textbf{Período de uso}: noviembre de 2025 a mayo de 2026.
    \item \textbf{Alcance}: Se utilizó como asistente de programación interactivo para acelerar la implementación de los componentes del sistema y la redacción de documentación técnica. Todo el código generado fue revisado, modificado cuando fue necesario y validado mediante las suites de tests unitarios y de extremo a extremo. Toda numerología verificable (tamaños de transacción, tasas de gas, órdenes de grupo) fue comprobada de forma independiente con cálculos en Python o ejecución directa de los crates.
    \item \textbf{Limitaciones}: La herramienta no tiene acceso a información posterior a su fecha de corte de entrenamiento. Las decisiones de diseño arquitectónico, la selección de algoritmos criptográficos y la definición de los objetivos del proyecto fueron responsabilidad exclusiva del autor.
  \end{itemize}

  \item \textbf{ChatGPT (OpenAI) --- familia GPT-4/5} y \textbf{GitHub Copilot}: utilizados puntualmente para sugerencias de autocompletado de código, búsquedas técnicas y comprobaciones de redacción de párrafos en español/inglés. Las salidas se incorporaron únicamente tras revisión manual.
\end{enumerate}

\subsection*{Aspectos no generados por IAG}

Los siguientes elementos fueron desarrollados íntegramente por el autor sin asistencia de IAG:

\begin{itemize}
  \item Definición del problema, objetivos y alcance del TFG.
  \item Selección de ML-DSA-65 frente a Falcon y SLH-DSA.
  \item Decisión de PoA frente a PoS para el consenso post-cuántico.
  \item Diseño de la arquitectura del sistema (capas, componentes, interfaces).
  \item Ejecución e interpretación de las simulaciones cuánticas con Qiskit.
  \item Ejecución, depuración y análisis de los tests E2E.
  \item Medición y análisis de los benchmarks de rendimiento.
  \item Revisión crítica y corrección de todo el contenido generado.
\end{itemize}

\subsection*{Declaración de responsabilidad}

El autor asume la plena responsabilidad sobre la corrección, originalidad y rigor académico de todo el contenido presentado en este trabajo, independientemente de las herramientas utilizadas durante su elaboración.
